% =====================================================================
% tesi/preambolo.tex - Preambolo della tesi
% =====================================================================
% Derivato da reports/preamble.tex e soggetto allo stesso vincolo, che va ripetuto
% perché è la ragione di quasi ogni scelta qui dentro: la TinyTeX di questa macchina
% è minimale e tlmgr rifiuta di installare pacchetti senza aggiornare prima se stesso.
% Ne segue che si usano soltanto pacchetti del nucleo, che la bibliografia non passa da
% BibTeX né da biblatex, e che dove un pacchetto comodo manca si definisce una macro
% di ricambio invece di aggiungere una dipendenza.
% ---------------------------------------------------------------------

\usepackage[T1]{fontenc}
\usepackage[utf8]{inputenc}
\usepackage{lmodern}
\usepackage[italian]{babel}
\usepackage[a4paper,margin=2.7cm]{geometry}
\usepackage{microtype}

% --- Tabelle e figure ---
\usepackage{booktabs}
\usepackage{tabularx}
\usepackage{longtable}
\usepackage{graphicx}
\usepackage[font=small,labelfont=bf]{caption}
\usepackage{xcolor}

% --- Matematica: serve per le temporizzazioni, i checksum e le quantizzazioni ---
\usepackage{amsmath}
\usepackage{amssymb}

% --- Codice sorgente ---
\usepackage{listings}
\definecolor{grigiochiaro}{gray}{0.96}
\definecolor{commento}{rgb}{0.35,0.45,0.35}
\lstset{
  basicstyle=\ttfamily\footnotesize,
  backgroundcolor=\color{grigiochiaro},
  commentstyle=\color{commento}\itshape,
  keywordstyle=\bfseries,
  breaklines=true,
  showstringspaces=false,
  frame=single,
  framerule=0pt,
  xleftmargin=0.5em,
  captionpos=b,
}
\lstdefinestyle{asm}{language={[x86masm]Assembler}}
\lstdefinestyle{c}{language=C}
\lstdefinestyle{py}{language=Python}

% --- Riferimenti interni ---
\usepackage[hidelinks,breaklinks=true]{hyperref}

% --- Indirizzi lunghi che non escono dal margine --------------------------------
% Per difetto \url non spezza un indirizzo, perché' un indirizzo interrotto può' non
% essere più' copiabile. In bibliografia questo produce righe che sbordano di centimetri,
% ed è esattamente ciò che accadeva alle voci con l'indirizzo di un thread di forum. La
% soluzione non è accorciare gli indirizzi, che devono restare verificabili, ma
% dichiarare dove l'interruzione è lecita: dopo una barra, un trattino, un punto e gli
% altri separatori che in un indirizzo non sono mai ambigui. Si aggiunge anche un po' di
% elasticità allo spazio interno, così' che la riga preferisca allungarsi invece di
% sbordare.
\usepackage{url}
\Urlmuskip=0mu plus 1mu\relax
\def\UrlBreaks{\do\/\do\-\do\_\do\.\do\?\do\=\do\&\do\%\do\+\do\:\do\~}
\def\UrlBigBreaks{\do\/\do\:}
% L'ultima difesa, per le righe che nemmeno così' rientrano: si concede al compositore di
% allargare gli spazi oltre il consueto prima di dichiarare una riga troppo piena.
\setlength{\emergencystretch}{3em}

\pagestyle{plain}

% --- Macro di ricambio per pacchetti non garantiti ---
% Lo spazio fine prima di un'unità di misura, che siunitx darebbe.
\providecommand{\SI}[2]{#1\,#2}
% addcontentsline esiste nel nucleo, ma il nome usato nel maestro è volutamente
% diverso per rendere evidente che si tratta di una voce aggiunta a mano all'indice.
\newcommand{\addcontentslineifsupported}[3]{\addcontentsline{#1}{#2}{#3}}

% --- Macro del progetto ---
% Un valore esadecimale, scritto sempre nello stesso modo in tutto il documento.
\newcommand{\hx}[1]{\texttt{0x#1}}
% Un valore binario.
\newcommand{\bin}[1]{\texttt{#1}\textsubscript{2}}
% Il livello di affidabilità di una fonte secondo la gerarchia di SOURCES.md.
\newcommand{\liv}[1]{\textsuperscript{L#1}}
% Un termine che entra nel glossario: si marca in corsivo alla prima occorrenza.
\newcommand{\term}[1]{\emph{#1}}
% Un percorso di file del progetto.
\newcommand{\file}[1]{\texttt{#1}}

% --- Riquadro per le affermazioni non verificate -------------------------------
% Il progetto impone di non presentare come fatto ciò che è inferito. In un documento
% lungo quella disciplina si perde se non ha una forma visibile, quindi le affermazioni
% da verificare vanno dentro questo ambiente e non nel corpo del testo.
% Deliberatamente costruito su quote e non su tabularx: tabularx analizza il proprio
% corpo per misurare le colonne, e non tollera di essere aperto e chiuso in due metà
% di un newenvironment. L'errore che ne segue è un runaway argument difficile da
% leggere, e un riquadro non vale una dipendenza fragile.
\newenvironment{daverificare}{%
  \par\medskip\begin{quote}\small\noindent
  \textcolor{gray}{\rule{2pt}{0.85em}}\ \textbf{Da verificare.}\ }{%
  \end{quote}\medskip}
